% !TEX root = ../notes_3dprinting.tex
% Chapter 10: Troubleshooting -- INTENTIONALLY EMPTY
% Grows from real failures on this machine. Do not pre-fill it with a
% generic catalogue; the outline below is a place to hang them, not a plan
% to execute.
%%%%%%%%%%%%%%%%%%%%%%%%%%%%%%%%%%%%%%%%%%%%%%%%%%%%%%%%%%%%%%%%%%%%%%

\index{troubleshooting}
\chapter{Troubleshooting}
\printchapterversion{10_troubleshooting}
\newthought{Filled by Failure.}

\newthought{Empty on purpose.} This chapter grows from real failures on this
machine rather than from a list of everything that could go wrong elsewhere.
When a print fails and you work out why, write it down here.

% Planned sections, uncomment and fill as failures actually occur:
% \section{Diagnose by reading the layer it failed on}
% \section{Filament snapped or ran out: the AMS case and the external case}
%   -> numbered procedure, the one place enumerate is justified
% \section{Clogs and no extrusion}
% \section{First layer failures}
% \section{Spaghetti, layer shift, aborted prints}
% \section{AMS faults}
% \section{Triage table: symptom, first check, fix}
